\documentclass[11pt]{article}
\usepackage[margin=1in]{geometry}
\usepackage{amsmath,amssymb,amsthm,mathtools}
\usepackage{graphicx}
\usepackage[hidelinks]{hyperref}
\usepackage{microtype}

\newtheorem{theorem}{Theorem}
\newtheorem{lemma}[theorem]{Lemma}
\newtheorem{corollary}[theorem]{Corollary}
\newcommand{\Z}{\mathbb Z}
\newcommand{\Q}{\mathbb Q}
\newcommand{\C}{\mathbb C}
\newcommand{\R}{\mathbb R}
\newcommand{\supp}{\operatorname{supp}}
\newcommand{\M}{M_{\mathcal D}}

\title{The five-digit irrational-mask example is never spectral}
\author{Candidate full negative answer to Example 5.2 of arXiv:2209.05619}
\date{11 August 2026}

\begin{document}
\maketitle

\begin{center}
\fbox{\parbox{0.92\linewidth}{\textbf{Status.}
Candidate full proof, likely valid, pending independent expert review.
In fact, the proof rules out even an infinite orthogonal family of
exponentials for every contraction ratio.}}
\end{center}

\section{Source problem}

An, He, and Lai end \emph{Classification of spectral self-similar measures
with four-digit elements} (arXiv:2209.05619) with the first genuinely
irrational five-digit test case \cite[Example 5.2]{AHL}.  For
\[
 \mathcal D=\{0,1,3,5,6\},\qquad 0<\rho<1,
\]
let $\mu_{\rho,\mathcal D}$ be the equal-weight self-similar measure for
the maps $x\mapsto\rho(x+d)$, $d\in\mathcal D$.  The source records that
all real zeros of the mask are irrational and states:

\begin{quote}
``We believe that none of such $\mu_{\rho,\mathcal D}$ can be a spectral
measure, but a more intensive study will be needed in this case.''
\end{quote}

\begin{center}
\includegraphics[width=0.97\linewidth]{figures/open_problem_crop.png}
\end{center}

We prove this belief for every $0<\rho<1$.  The argument also gives a
general obstruction whenever an integral mask has no rational zero phase.

\section{Proof intuition}

The Fourier transform is an infinite mask product.  Hence every difference
of two orthogonal frequencies has the form
\[
 q^k(\theta+n),\qquad q=\rho^{-1}>1,quad k\geq1,quad n\in\Z,
\]
where $\theta$ is one of finitely many mask-zero phases.  Infinite Ramsey
theory gives an infinite subfamily for which the same $\theta$ works for
every pair.

The phases in the source example are not merely irrational: their
exponentials are algebraic, so Gelfond--Schneider makes the phases
transcendental.  If $q$ is algebraic, coefficient comparison then forces
pairwise differences of infinitely many powers of $q$ to be powers of $q$,
which is impossible by lacunarity.

If $q$ is transcendental, one difference relation writes
$\theta=R(q)$ for a nonconstant rational function $R\in\Q(X)$.  Every other
numerical difference identity then lifts, by transcendence of $q$, to an
identity in $\Q(X)$.  A degree-at-infinity lemma shows that the set
$\{X^k(R+n)\}$ cannot contain an infinite one-sided difference clique.
This eliminates the remaining case.

\section{Preliminaries}

For a finite integral digit set $\mathcal D$, write
\[
 M_{\mathcal D}(\xi)=\frac1{\#\mathcal D}
   \sum_{d\in\mathcal D}e^{-2\pi i d\xi},
 \qquad
 \Theta=\mathcal Z(M_{\mathcal D})\cap[0,1).
\]
The Fourier transform of the associated equal-weight self-similar measure is
\begin{equation}\label{eq:product}
 \widehat\mu_{\rho,\mathcal D}(\xi)
   =\prod_{j=1}^{\infty}M_{\mathcal D}(\rho^j\xi).
\end{equation}
For fixed $\xi$, the product is nonzero if none of its factors is zero,
because $M_{\mathcal D}(t)=1+O(t)$ at the origin and
$\sum_j|\rho^j\xi|<\infty$.  Therefore, with $q=\rho^{-1}$,
\begin{equation}\label{eq:zeros}
 \mathcal Z(\widehat\mu_{\rho,\mathcal D})
   =\bigcup_{\theta\in\Theta}\bigcup_{k\geq1}
       q^k(\theta+\Z).
\end{equation}

If $E(\Lambda)=\{e^{2\pi i\lambda x}:\lambda\in\Lambda\}$ is
orthogonal in $L^2(\mu)$, then
\begin{equation}\label{eq:bizero}
 (\Lambda-\Lambda)\setminus\{0\}
    \subseteq\mathcal Z(\widehat\mu).
\end{equation}

\begin{lemma}[Transcendence of irrational integral-mask phases]
\label{lem:phase-trans}
Let $\mathcal D\subset\Z$ be finite.  If
$\theta\in\mathcal Z(M_{\mathcal D})$ is irrational, then $\theta$ is
transcendental.
\end{lemma}

\begin{proof}
The number $z=e^{-2\pi i\theta}$ is a zero of the integral polynomial
$\sum_{d\in\mathcal D}z^d$ (after multiplying by a power of $z$ if
necessary), so $z$ is algebraic.  If $\theta$ were algebraic irrational,
then Gelfond--Schneider, applied to the algebraic base $-1$ and algebraic
irrational exponent $-2\theta$, would say that every value of
$(-1)^{-2\theta}$ is transcendental.  Taking $\log(-1)=\pi i$ gives the
value $e^{-2\pi i\theta}=z$, a contradiction.
\end{proof}

The next lemma is the formal obstruction needed only when $q$ is
transcendental.  For a nonzero rational function $F$, let
$\deg_\infty F$ be numerator degree minus denominator degree in reduced form.

\begin{lemma}[Rational-function difference-clique obstruction]
\label{lem:formal}
Let $R\in\Q(X)\setminus\Q$ and
\[
 \mathcal S_R=\{X^k(R+n):k\geq1,\ n\in\Z\}\subset\Q(X).
\]
There is no infinite sequence of distinct elements
 $f_0,f_1,\ldots\in\mathcal S_R$ such that
\begin{equation}\label{eq:formalclique}
 f_j-f_i\in\mathcal S_R\qquad(0\leq i<j).
\end{equation}
\end{lemma}

\begin{proof}
We use only the behavior at $X=\infty$.

First suppose $R$ has a pole there, of order $d>0$.  Every
$X^k(R+n)$ has degree $k+d$.  An infinite fixed-$k$ subfamily is
impossible: the difference of two members is $(n'-n)X^k$, while equality
with $X^c(R+\ell)$ forces $c=k-d$ and
$R+\ell=(n'-n)X^d$; three distinct $n$'s would force two different
coefficients for the same $R$.  Thus, after passing to a subsequence, its
levels $k_j$ tend to infinity.  Fix $f_0$.  For all large $j$ the degree of
$f_j-f_0$ is $k_j+d$, so \eqref{eq:formalclique} has level $k_j$ and
\[
 X^{k_j}(R+n_j)-f_0=X^{k_j}(R+\ell_j).
\]
Hence $f_0=(n_j-\ell_j)X^{k_j}$, which cannot hold for two distinct
unbounded $k_j$.

Now suppose $L=R(\infty)\in\Q$.  Put $H=R-L$.  A level $(k,n)$ is
\emph{normal} if $L+n\ne0$; it then has degree $k$.  If $-L\in\Z$, there
is also one exceptional choice $n_*=-L$, for which
$X^k(R+n_*)=X^kH$.  Write $\deg_\infty H=-d<0$.

Assume first that infinitely many $f_j$ are normal.  As in the preceding
paragraph, an infinite fixed-level subfamily is impossible (differences are
integer multiples of a monomial, while a possible exceptional
representation fixes that multiple).  Pass to normal levels $k_j\to\infty$
and fix $f_0$.  If $f_j-f_0=X^{c_j}(R+\ell_j)$ is normal, degrees give
$c_j=k_j$ and then $f_0=(n_j-\ell_j)X^{k_j}$, again impossible twice.  If it
is exceptional, then $c_j=k_j+d$ and
\begin{equation}\label{eq:Fm}
 f_0=X^{k_j}\bigl[(L+n_j)+H-X^dH\bigr].
\end{equation}
As $n$ varies, the bracketed rational functions differ only by integer
constants.  All but at most one have degree zero; the remaining one has a
single fixed finite degree.  Their degrees therefore cannot equal
$\deg_\infty f_0-k_j\to-\infty$, contradicting \eqref{eq:Fm}.

It remains to consider infinitely many exceptional terms $X^{k_j}H$;
their levels are necessarily distinct and unbounded.  Fix $k_0$.  If
\[
 H(X^{k_j}-X^{k_0})=X^c(R+\ell)
\]
has an exceptional right side, cancellation of $H$ would give
$X^{k_j}-X^{k_0}=X^c$, an impossible Laurent-polynomial identity.  For a
normal right side, degrees give $c=k_j-d$, and rearrangement yields
\begin{equation}\label{eq:constant}
 H\bigl(X^d-1-X^{k_0-k_j+d}\bigr)=L+\ell\in\Q.
\end{equation}
Three distinct values of $k_j$ in \eqref{eq:constant}, after subtracting
pairs, would make the ratio of two distinct binomials
\[
 \frac{X^a-X^b}{X^a-X^c}
\]
a rational constant, which it is not.  Thus this case is finite as well.
All possibilities contradict infinitude.
\end{proof}

\section{The general obstruction}

\begin{theorem}\label{thm:general}
Let $\mathcal D\subset\Z$ be finite with at least two elements, and suppose
every phase in $\Theta=\mathcal Z(M_{\mathcal D})\cap[0,1)$ is irrational.
Then, for every $0<\rho<1$, the space
$L^2(\mu_{\rho,\mathcal D})$ contains no infinite orthogonal family of
exponentials.  In particular, $\mu_{\rho,\mathcal D}$ is not spectral.
\end{theorem}

\begin{proof}
Suppose $\Lambda$ is an infinite orthogonal set.  Color each unordered pair
of distinct elements of $\Lambda$ by one phase $\theta\in\Theta$ occurring
in a representation \eqref{eq:zeros} of the absolute difference (choose one
deterministically if there is more than one).  The set of colors is finite.
By the infinite Ramsey theorem, there is an infinite monochromatic subset.
It has an infinite strictly monotone subsequence; reverse its sign if needed
and write it as
\[
 \gamma_0<\gamma_1<\gamma_2<\cdots.
\]
For one fixed $\theta\in\Theta$, every positive difference is therefore
\begin{equation}\label{eq:monochrome}
 \gamma_j-\gamma_i=q^{k_{ij}}(\theta+n_{ij}),
 \qquad i<j,quad k_{ij}\geq1,quad n_{ij}\in\Z.
\end{equation}
By Lemma \ref{lem:phase-trans}, $\theta$ is transcendental.

Translate by $\gamma_0$ and set
\[
 d_i=\gamma_i-\gamma_0=q^{k_i}(\theta+n_i)\qquad(i\geq1).
\]

If $q$ is algebraic, compare $d_j-d_i$ with
\eqref{eq:monochrome}.  Since $\theta$ is transcendental over the algebraic
numbers, its coefficient must vanish:
\begin{equation}\label{eq:powerdifference}
 q^{k_j}-q^{k_i}=q^{k_{ij}}\qquad(i<j).
\end{equation}
Thus $k_j>k_i$.  Fix $i=1$ and let $j\to\infty$.  Dividing
\eqref{eq:powerdifference} by $q^{k_j}$ gives
\[
 1-q^{k_1-k_j}=q^{k_{1j}-k_j}.
\]
The right side is a power of $q$ strictly below $1$, hence at most $q^{-1}$,
whereas the left side tends to $1$ from below.  This is impossible.

It remains that $q$ is transcendental.  Use $d_2-d_1$ once.  On collecting
the coefficient of $\theta$, we obtain
\[
 \bigl(q^{k_2}-q^{k_1}-q^{k_{12}}\bigr)\theta=B(q)
\]
for a Laurent polynomial $B\in\Z[X,X^{-1}]$.  The coefficient polynomial is
not identically zero (its value at $X=1$ is $-1$), and it cannot vanish at
the transcendental number $q$.  Hence
\[
 \theta=R(q)
\]
for some $R\in\Q(X)$.  This $R$ is nonconstant because $\theta$ is
irrational.

Replace $q$ formally by $X$ and define
\[
 f_i=X^{k_i}(R+n_i)\in\mathcal S_R.
\]
Every equality \eqref{eq:monochrome}, after substituting
$\theta=R(q)$, is an equality of rational functions evaluated at the
transcendental number $q$.  A nonzero rational function over $\Q$ cannot
vanish at $q$, so each equality lifts identically to $\Q(X)$.  Consequently
$f_j-f_i\in\mathcal S_R$ for all $i<j$, contradicting Lemma
\ref{lem:formal}.

Thus no infinite orthogonal exponential set exists.  The self-similar
attractor is infinite when $\#\mathcal D\geq2$, so $L^2(\mu)$ is
infinite-dimensional; a finite orthogonal family cannot be an orthonormal
basis.  Hence the measure is not spectral.
\end{proof}

\begin{corollary}[The source's five-digit example]\label{cor:source}
For $\mathcal D=\{0,1,3,5,6\}$ and every $0<\rho<1$, the measure
$\mu_{\rho,\mathcal D}$ is not spectral.  More strongly, it admits no
infinite orthogonal set of exponentials.
\end{corollary}

\begin{proof}
The source paper records
$\mathcal Z(M_{\mathcal D})=\Theta+\Z$ with
$\Theta\subset\R\setminus\Q$.  Apply Theorem \ref{thm:general}.
\end{proof}

\section{Verification, scope, and novelty}

\paragraph{Proof stress tests.}
The Fourier-product zero equality uses absolute convergence near the origin,
so it introduces no spurious zeros.  Ramsey colors only by the finite phase,
not by the unbounded scale.  The algebraic-$q$ argument covers all algebraic
contractions, not merely reciprocal integers.  In the transcendental case,
evaluation at $q$ is injective on $\Q(X)$; the rational-function lemma treats
separately poles at infinity, ordinary finite limits, and the sole possible
exceptional integer translate.  Zero mask phases and repeated
representations cause no issue because the source phases are irrational and
the coloring chooses one representation.

\paragraph{Scope.}
Corollary \ref{cor:source} completely proves the explicit belief in Example
5.2, for all contraction ratios.  Theorem \ref{thm:general} also covers every
finite integral digit set whose mask has no rational zero phase.  It does not
settle the full modified {\L}aba--Wang conjecture when rational and irrational
mask phases coexist.

\paragraph{Bounded novelty check.}
On 11 August 2026 we searched the run's cheap registry/problem indexes, the
exact digit set, the exact source sentence, the source title and citations,
and close phrases involving irrational mask zeros and infinite orthogonal
exponentials.  We found the source paper and later works citing its four-digit
classification, but no later resolution of Example 5.2 or statement matching
Theorem \ref{thm:general}.  The ingredients (infinite Ramsey and
Gelfond--Schneider) are classical; their combination with the
rational-function difference-clique lemma appears new.  Novelty confidence is
moderate, while proof confidence is high.

\paragraph{Human-review recommendation.}
Verify the degree-at-infinity case split in Lemma \ref{lem:formal}, especially
the exceptional translate $n=-R(\infty)$; then check the lifting of numerical
identities to $\Q(X)$ when $q$ is transcendental.  The included symbolic
script is only a sanity check and is not used as proof.

\begin{thebibliography}{9}
\bibitem{AHL}
L.-X. An, X. He, and C.-K. Lai,
\emph{Classification of spectral self-similar measures with four-digit
elements}, Asian J. Math. 27 (2023), 467--492; arXiv:2209.05619.

\bibitem{Baker}
A. Baker,
\emph{Transcendental Number Theory},
Cambridge University Press, 1975.

\bibitem{Ramsey}
F. P. Ramsey,
\emph{On a problem of formal logic},
Proc. London Math. Soc. (2) 30 (1930), 264--286.
\end{thebibliography}

\end{document}
